%=====================================================================================
\section{Introduction}\label{sec:intro}
%=====================================================================================
The vehicle routing problem with time windows (VRPTW) asks for minimum-cost routes that
serve every customer within its time window. It is NP-hard, and Adaptive Large
Neighbourhood Search (ALNS)~\cite{ropke2006alns,pisinger2007general} is one of the most widely
used metaheuristics for it~\cite{mara2022survey}. At every iteration, a \emph{destroy} operator removes $q$
customers from the incumbent, a \emph{repair} operator re-inserts them, and a
simulated-annealing (SA) criterion decides acceptance; a roulette wheel (RW) biases
operator selection towards recently successful operators. A
meta-analysis~\cite{turkes2021meta}, however, found that this adaptive layer improves
solution quality by only about 0.14\% on average. This makes ALNS a natural target for
learned \emph{algorithm control}~\cite{karimi2022survey}, also known as parameter
control~\cite{eiben1999control} or dynamic algorithm configuration
(DAC)~\cite{biedenkapp2020dac,adriaensen2022dac}. Recent works train deep
reinforcement learning (RL) agents to control ALNS at every
iteration~\cite{kallestad2023,reijnen2024,johnn2024grlos,wang2025ppoalns} and report
gains over RW.

These agents share three limitations. First, to keep the Markov decision process (MDP)
tractable, they are trained on short episodes (20 decisions in
GRLOS~\cite{johnn2024grlos}, 100 iterations in DR-ALNS~\cite{reijnen2024}) or on runs
shorter than at deployment (1{,}000 vs.\ up to 10{,}000 iterations in
DRLH~\cite{kallestad2023}), and Johnn et al.~\cite{johnn2024grlos} report degradation
outside the training horizon. Second, most of them fix the destroy scale and the temperature, although both
strongly affect ALNS~\cite{ropke2006alns,santini2018}; DR-ALNS~\cite{reijnen2024} and
recent variants~\cite{cherrered2026roadef,xie2025pod} also learn them. Third, to our knowledge, none of them can fall
back to the RW rule as one of its actions.

More fundamentally, it is rarely checked \emph{why} the agent wins: does it exploit the
search state, or has it merely found a better static configuration? Control subsumes
configuration, so a controller can win without being adaptive; DAC studies therefore
compare learned policies with the best static
configuration~\cite{biedenkapp2020dac,eimer2021dacbench}. In RL-controlled ALNS,
Reijnen et al.~\cite{reijnen2024} compare with an ALNS tuned by SMAC, but whether the
learned policies actually depend on the state, and how they vary across training runs,
has not been reported.

We study this question with \sgalns{} (Segment-level Graph ALNS), a controller designed
to be favourable to RL. Following the DAC formulation, in which a policy
reconfigures the algorithm at fixed intervals~\cite{biedenkapp2020dac}, one decision
configures the next $K$ iterations, so that proximal policy optimisation (PPO) can be
trained on complete ALNS runs. Each decision jointly selects the operators, the destroy
scale and the temperature from a graph state with dynamic time-window features, and the
action space contains the RW rule. Our contributions are:
\begin{itemize}
\item \textbf{A segment-level MDP for ALNS} on the VRPTW, used as a deliberately
favourable test bed rather than as a new principle: it adopts interval-based decisions
from DAC~\cite{biedenkapp2020dac} and joint control of operators, destroy scale and
temperature from DR-ALNS~\cite{reijnen2024}, and adds full-trajectory training, a
dynamic time-window graph state and an \emph{expert-inclusive} action space.
\item \textbf{An evaluation protocol with six controls} for RL-controlled ALNS. Static
baselines and Friedman--Nemenyi tests are established practice in DAC and algorithm
benchmarking~\cite{eimer2021dacbench,demsar2006}; we combine them with a policy-collapse
diagnostic, training-seed replication, tuned static \emph{and} tuned \Rw{} baselines
with a tuning budget smaller than one training run, a dynamic-headroom test, a
time-matched comparison, and transfer across budgets and instance sizes up to $n=400$,
all evaluated with hierarchical two-sided tests and equivalence tests.
\item \textbf{Evidence that learned ALNS control on the VRPTW is, in our setting,
learned configuration}: \tCollapsed{} of the \tNpol{} trained policies choose a
constant action on the instances and budgets we examined and ignore counterfactual
changes of their state, while the exceptions are worse or, in one training seed, a
time-dependent schedule no better than a tuned constant; seed-to-seed differences dominate, tuned
configurations are equivalent in fleet size, time-dependent schedules offer small
headroom that a direct search finds but the policies do not exploit, and transfer
follows the training distribution rather than any adaptation.
\end{itemize}
Table~\ref{tab:controls} summarises the six controls and their outcome; \Sect~\ref{sec:results}
details them.
{\widetables
\begin{table}[t]
\centering
\caption{The six controls at a glance}\label{tab:controls}
\setlength{\tabcolsep}{4pt}\small
\begin{adjustbox}{max width=\linewidth}
\begin{tabular}{lp{9.2cm}}
\toprule
Control & Result\\ \midrule
(1) Collapse & 17 of 20 policies use one action in all decisions (Fig.~\ref{fig:behaviour}); \qProbeChg{}\% of the counterfactual queries change the action of a collapsed policy (Table~\ref{tab:probe}) \\
(2) Training seeds & CNV 419.3$\pm$0.9 over 5 seeds; 4 different actions; 2 of 5 use significantly fewer vehicles than \Rw{} \\
(3) Tuned baselines & vs.\ tuned static: $\Delta$CNV $-$0.6 [$-$2.0, +0.8], $p$=0.581 (Table~\ref{tab:stats}) \\
(4) Headroom & best two-phase schedule \vdhDyn{} vs.\ best constant \vdhConst{} ($p$=\vdhP{}); per-class oracle saves \qOrGain{} vehicles over the best constant \\
(5) Equal run time & vs.\ tuned static at $T$=10k: $\Delta$CNV +0.2 [$-$1.2, +1.6], $p$=1.000 \\
(6a) Budget & vs.\ tuned static, $T$=1k: $\Delta$CNV $-$0.7 ($p$=0.518); $T$=20k: $\Delta$CNV +1.0 ($p$=0.336) \\
(6b) Size $n$=400 & vs.\ \Rw{}: \sgalns{} $\Delta$TD +2.37\% ($p<0.001$), \sgc{} $-$1.06\% ($p<0.001$) on equal-fleet instances \\
\bottomrule
\end{tabular}
\end{adjustbox}
\par\vspace{2pt}\parbox{\linewidth}{\footnotesize Solomon, $T=5{,}000$ unless stated; $\Delta$ = \sgalns{} minus baseline, two-sided tests; details in Tables~\ref{tab:compare} and~\ref{tab:stats}.}
\end{table}
}
